\section{Strict higher categories and transformation homotopy}

Throughout this section, morphisms of strict $n$-categories are strict functors. Transformation homotopy is the equivalence relation generated by normal lax transformations, or by pseudonatural transformations in dimension two. For arbitrary $n$, the lax convention is the Gray-cylinder convention, with its strict unit and composition identities. The pseudo subcase uses invertible structural cells. Let $\mathcal H_n$ denote the resulting quotient category. Write $B_nX=|N_nX|$ for Street's nerve, and regard ordinary groupoids as strict $n$-categories with identity higher cells.

\begin{claim}{E.1}{Refuted}{The proposed classification by two-types and two-groupoids}
The strict $2$-groupoid $\mathbf B^2C_2$ with one object, one $1$-arrow, and $2$-automorphism group $C_2$ is not local in either transformation convention. For $X=\mathbf BC_2$, the canonical comparison has a singleton domain and a two-element codomain:
\[
 \Hom_{\mathcal H_2}(\mathbf BC_2,\mathbf B^2C_2)
 \longrightarrow [BC_2,K(C_2,2)]. \tag{E.1.1}
\]
Consequently the statement that all strict $2$-groupoids are local is false.
\end{claim}
\begin{proof}
A strict functor sends the unique object to the unique object and every $1$-arrow to the unique $1$-arrow. There are no nonidentity source $2$-cells, so exactly one strict functor exists. Taking any equivalence relation on this singleton leaves a singleton.

The Street nerve of $\mathbf B^2A$ is the standard $K(A,2)$: its simplices are normalized $A$-valued $2$-cocycles on simplices, and the tetrahedral identity is the cocycle identity. Equivalently its Moore complex has $A$ in degree two and zero elsewhere. For $A=C_2$, a normalized group $2$-cocycle on $C_2$ is determined by its value on the pair $(1,1)$; either value satisfies the cocycle identity. A normalized $1$-cochain has zero coboundary on this pair. Hence $H^2(C_2;C_2)=C_2$, proving the assertion about the codomain. This is an unpointed computation as well, since the target is simply connected. The program \texttt{code/e\_strict\_two\_groupoid.py} exhausts the relevant cochains and verifies both compositions and interchange in the finite strict $2$-groupoid.
\end{proof}

\begin{claim}{E.2}{Proved}{The groupoid reflection in every strict dimension}
For $1\leq n\leq\infty$, let $\Pi X$ be the groupoid with generators the $1$-cells of $X$, relations for their strict compositions, and the relation $f=g$ for every $2$-cell $f\Rightarrow g$. The strict functor $\eta_X:X\to\Pi X$ induces the fundamental groupoid of $B_nX$. It gives a reflection
\[
 \Pi:\mathcal H_n\rightleftarrows\Ho(\Grpd):j.
\]
For each groupoid $G$, its canonical comparison is bijective:
\[
 \Hom_{\mathcal H_n}(X,G)\cong[B_nX,BG]. \tag{E.2.1}
\]
\end{claim}
\begin{proof}
The edges of $N_nX$ are the $1$-cells. Its triangular simplices record a $2$-cell from one composite to another edge. Fundamental groupoidification imposes exactly these triangular relations. Degenerate-edge triangles include every relation given by a $2$-cell, and triangles with identity $2$-cell include the composition relations. Higher simplices impose no additional relation on the fundamental groupoid. This proves the specified identification with $\Pi_1B_nX$.

A strict functor $X\to G$ must send every higher cell to an identity, and factors uniquely through $\eta_X$. In a lax transformation with target $G$, every structural $2$-cell is an identity; its component $1$-arrows are invertible, and its equations are precisely the ordinary naturality equations. Thus transformations are natural isomorphisms after passage to $\Pi X$, giving the reflection on the quotient Hom-sets. Maps of spaces to a groupoid nerve are exactly fundamental-groupoid functors modulo natural isomorphism, as checked in B.1. This proves (E.2.1), with its canonical map.
\end{proof}

\begin{claim}{E.3}{Proved}{The quotient localizes to ordinary homotopy types}
The functor $\mathcal H_n\to\Hot_\U$ induced by $B_n$ is localization at Street-nerve weak equivalences, for every $1\leq n\leq\infty$.
\end{claim}
\begin{proof}
Theorem 5.6 of \citet{Gagna2018}, opened with its proof, identifies $n\text{-}\Cat[N_n^{-1}(W_\infty)^{-1}]$ with $\Hot$. Its hypotheses apply to small strict globular $n$-categories and the Street nerve used here, including $n=\infty$. Here is the relevant comparison underlying that theorem. Put $Q=\operatorname{Sd}^2 i_!i^*$, with $i$ the inclusion of injective simplex maps. There is a natural weak equivalence $QZ\to Z$, and $QZ$ is the nerve of a poset. The poset-unit comparison $QZ\to N_nc_nQZ$ is a weak equivalence. Its proof is Gagna's explicit homotopy in Theorem 6.9. The triangle identities now show that $c_nQ$ is an inverse at localization: the counit map $c_nQN_nX\to X$ becomes a weak equivalence after applying $N_n$. This checks the precise localization statement needed, rather than assuming a transferred model structure.

A Gray lax or oplax transformation induces a homotopy of Street nerves. The explicit map is obtained from the Alexander--Whitney diagonal on the simplex chains:
\[
 [i_0,\ldots,i_p]\longmapsto
 \sum_{a=0}^p[i_0,\ldots,i_a]\otimes[i_a,\ldots,i_p].
\]
Its differential commutes with the sum because the two terms deleting a shared vertex cancel with signs $(-1)^a$ and $(-1)^{a-1}$. All summands are positive; the augmentation is preserved. Naturality under a degeneracy follows because each summand has a repeated entry in one tensor factor. More explicitly, let $\mathcal O_m$ be the $m$th oriental, let $D_1$ be the directed interval, and write $\delta_m=\nu\mathrm{AW}_m:\mathcal O_m\to\mathcal O_m\otimes\mathcal O_m$. For an $m$-simplex $F:\mathcal O_m\to X$ and $\theta:[m]\to[1]$, the formula
\[
 J_X(F,\theta)=(F\otimes\mathcal O_\theta)\delta_m:
 \mathcal O_m\longrightarrow X\otimes D_1
\]
is natural in simplex maps by the checked chain formula. A transformation represented by $P:X\otimes D_1\to Y$ sends this simplex to $P J_X(F,\theta)$ and therefore gives $N_nX\times\Delta^1\to N_nY$. Its restrictions at $0,1$ are the two given strict functors. Reversing the tensor order gives the opposite lax convention with the same endpoints. This checks the nerve-homotopy lemma in the attached OpenAI manuscript, Section 2, and agrees with \citet[\S3.3 and Proposition 3.4]{AraTheoremB}. The claimed (OpenAI, 2026, unrefereed) transferred model structure is not required in this proof.

Consequently the localization functor on $n\text{-}\Cat$ identifies the transformation relation. Factoring it through $\mathcal H_n$ and applying the universal property of localization proves the claim. The pseudo relation is a subrelation consisting of transformations with the same nerve-homotopy property, so the identical localization argument applies.
\end{proof}

\begin{claim}{E.4}{Proved}{Descent of transformations along the R2 quotients}
Let $C$ be a strict $n$-category with at most $\mu$ cells in all dimensions, where $\mu$ is infinite. Let $p:G\twoheadrightarrow\Gamma$ be a group quotient such that every homomorphism from $G$ to a group of cardinality at most $\mu$ kills $\ker p$. Restriction along $\mathbf Bp$ is an isomorphism on the categories of strict functors and normal lax transformations from $\mathbf B\Gamma$ and $\mathbf BG$ to $C$. It induces a bijection on transformation-homotopy classes, including the pseudo convention.
\end{claim}
\begin{proof}
A strict functor $\mathbf BG\to C$ consists of an object $x$ and a homomorphism to the group of strictly invertible $1$-endomorphisms of $x$. That group has size at most $\mu$, so the functor factors uniquely through $p$.

To check transformations, let their object component be a $1$-cell $u:x\to y$. Define $T_u$ to consist of triples $(f,g,\alpha)$, where $f:x\to x$ and $g:y\to y$ are strictly invertible $1$-cells and $\alpha:uf\Rightarrow gu$ is an invertible $2$-cell. Its multiplication is
\[
 (f,g,\alpha)(f',g',\alpha')
 =\bigl(ff',gg',(g\ast\alpha')\circ(\alpha\ast f')\bigr).
\]
The strict interchange, associativity, and unit identities verify the group law. Its identity is $(1,1,1_u)$ and its inverse is given by whiskering $\alpha^{-1}$ with $g^{-1}$ and $f^{-1}$. There are at most $\mu$ triples.

For a transformation on $\mathbf BG$, its cell at an element $h\in G$ is invertible even for the lax convention: the transformation equations at $hh^{-1}=1$ and $h^{-1}h=1$ provide its two inverses after whiskering. The transformation's composition and unit equations say exactly that these triples define a homomorphism $G\to T_u$. This homomorphism kills $\ker p$ and hence descends uniquely. For a group as source there are only identity higher cells; the normal Gray transformation identities make their components identities, so no further data remain in dimensions above two. The descended transformations preserve vertical composition because their restrictions do, and restriction is faithful on their components. This proves the actual isomorphism before taking components.
\end{proof}

\begin{claim}{E.5}{Proved}{Invariant small ind-objects in every strict dimension}
For $1\leq n\leq\infty$, the Street-nerve invariant $\U$-small ind-objects of $\mathcal H_n$ are exactly the ind-objects of ordinary groupoids. For an invariant diagram $(K_i)_{i\in I}$ the canonical morphism
\[
 (K_i)_{i\in I}\longrightarrow(\Pi K_i)_{i\in I}
\]
is an isomorphism in $\Ind_\U(\mathcal H_n)$. The same assertion about canonical presheaves holds for arbitrary $\U$-small colimits of representables.
\end{claim}
\begin{proof}
Choose one infinite $\mu\in\U$ bounding the total cells of every $K_i$; this includes countably many dimensions when $n=\infty$. By E.3, the Yoneda descent calculation of B.1 gives $F(Y)=\colim_i[Y,B_nK_i]$. Apply R2 to each component of $Y$, using this bound, to obtain $T=\coprod_cBG_c\to Y$ and the quotient $T\to\tau_1Y$. E.4 identifies the transformation-Hom sets from the two groupoid domains to every $K_i$. R2 also gives bijections $[Y,B_nK_i]\to[T,B_nK_i]$, since the fundamental groups of $B_nK_i$ have at most $\mu$ elements. The same commuting triangle as in B.1 proves $F(\tau_1Y)\to F(Y)$ bijective. E.2 identifies this map with precomposition by $\eta_X$ for $Y=B_nX$.

The ind-reflection argument checked in B.1 now gives both inverse identities for the displayed canonical morphism. The arbitrary-colimit conclusion follows from left Kan extension as in B.2. Conversely all groupoid representables are nerve invariant by (E.2.1), and any colimit of their bijections remains a bijection.
\end{proof}

\begin{claim}{E.6}{Proved}{Representable homotopy types and local objects}
For each $1\leq n\leq\infty$, a small space $Y$ has $X\mapsto[B_nX,Y]$ representable on $\mathcal H_n$ precisely when $Y$ is a $1$-type. The local objects are exactly those isomorphic in $\mathcal H_n$ to ordinary groupoids, equivalently those whose canonical map $K\to\Pi K$ is an isomorphism there. These conclusions apply to both requested transformation conventions in dimension two.
\end{claim}
\begin{proof}
A represented invariant presheaf is a constant-index case of E.5, so its representing object $K$ is isomorphic to the groupoid $\Pi K$. Alternatively, a representation of $[B_n-,Y]$ gives a representing map $B_nK\to Y$; E.3 and Yoneda make this map invertible in $\Hot$. It follows that $Y$ is a $1$-type. Every $1$-type is represented by its fundamental groupoid by E.2. For local objects apply this argument to $Y=B_nK$. Conversely an isomorphism to a groupoid transports the canonical comparison through a commuting square of bijections. The result settles the stated strict-$n$ question; the expected larger class of nerves of strict $n$-groupoids already contains the explicit counterexample E.1.
\end{proof}
